\usepackage[
    papersize={170mm,240mm},
    left=20mm,
    right=20mm,
    top=15mm,
    bottom=15mm,
    includeheadfoot=true,
    headheight=12.6pt
]{geometry}
\usepackage[LGR,T2A]{fontenc}
\usepackage[russian]{babel}
\usepackage{indentfirst}

%--- форматирование
\usepackage{pdflscape}
\usepackage{float}
\usepackage{multirow}
\usepackage{adjustbox}
\usepackage{titlesec}
\usepackage{enumitem}
\setlist[enumerate]{noitemsep, topsep=1mm}
\setlist[itemize]{noitemsep, topsep=1mm}
\usepackage{graphicx}
\usepackage{wrapfig}
\usepackage{amssymb, amsmath}
\usepackage{wasysym} % для символа в блоке потока терминов
\usepackage{pifont} % для цифр в окружности (для диаграмм)
\usepackage[nopatch=footnote]{microtype}

%--- переопределение формата заголовков
\titlespacing{\subsection}{0pt}{*2}{*1}
\titleformat{\subsection}
    {\normalfont\large\bfseries}
    {\thesubsection}
    {1em}
    {}

%--- оформление задач и примеров
\usepackage{tcolorbox}
\tcbuselibrary{theorems}
\tcbuselibrary{skins}
\tcbuselibrary{breakable}
\usepackage{answers}
\newtcbtheorem[number within=section]
    {problem}
    {Задача}
    {
        coltitle=black,
        enhanced,
        attach boxed title to top left={
            xshift=5mm,
            yshift=-1.3mm
        },
        boxed title style={
            colback=yellow!20,
            top=-0.5mm,
            bottom=-0.8mm,
            breakable,
            boxrule=0.5pt,
            colframe=red!60
        },
        colback=white,
        top=1mm,
        bottom=0mm,
        left=0.5mm,
        right=0.5mm,
        before skip=2mm,
        after skip=2mm,
        fontupper=\small,
        fonttitle=\small,
        breakable,
        boxrule=0.5pt,
        colframe=red!60
    }
    {prob}
\Newassociation{sol}{solution}{solutions}
\renewcommand{\solution}[1]{\ignorespaces}
\Opensolutionfile{solutions}[ans]
\newtcbtheorem[number within=section]
    {example}
    {Пример}
    {
        coltitle=black,
        enhanced,
        attach boxed title to top left={xshift=5mm, yshift=-1.3mm},
        boxed title style={
            colback=blue!5,
            top=-0.5mm,
            bottom=-0.8mm,
            boxrule=0.5pt,
            colframe=blue!60
        },
        colback=white,
        top=1mm,
        bottom=0mm,
        left=0.5mm,
        right=0.5mm,
        before skip=2mm,
        after skip=2mm,
        fontupper=\small,
        fonttitle=\small,
        breakable,
        boxrule=0.5pt,
        colframe=blue!60
    }
    {ex}

\newenvironment{condition}{\par\itshape}{\par} % блок условий

%--- формат потока терминов
\newtcolorbox{termlist}{
    toptitle=1pt,
    bottomtitle=0.5pt,
    title={$\photon$ \textbf{Поток терминов} $\photon \photon \photon$},
    colframe=green_sea,
    colback=white,
    boxrule=0.5pt,
    top=1pt,
    bottom=1pt,
    left=0.5mm,
    right=0.5mm,
    before skip=2mm,
    after skip=2mm,
    fontupper=\small,
    fonttitle=\small,
    colbacktitle=white,
    coltitle=black
}

%--- настройка формата подписей к рисункам
\usepackage{caption}
\captionsetup{font={it,small}}

%--- настройка колонтитулов
\usepackage{fancyhdr}
\pagestyle{fancy}
\fancyhf{}
\fancyhead[LE,RO]{\thepage}
\fancyhead[RE]{\leftmark}
\fancyhead[LO]{\rightmark}

%--- настройка таблиц
\usepackage{xcolor}
\usepackage{colortbl}
\usepackage{tabularx}
\usepackage{ragged2e}
\usepackage{array}
\newcolumntype{Y}{>{\RaggedRight\hspace{0pt}\arraybackslash}X}
\newcolumntype{Z}{>{\centering\arraybackslash}X}
\setlength{\tabcolsep}{6pt}
\renewcommand{\arraystretch}{1.3}

%--- форматирование оглавления
\usepackage{tocloft}
\setlength{\cftbeforechapskip}{5pt}
\setlength{\cftsecindent}{1.5em}
\setlength{\cftsecnumwidth}{2.3em}

%--- греческие слова в тексте
\newcommand{\grk}[1]{{\fontencoding{LGR}\selectfont #1}}

%--- предметный указатель
\usepackage{makeidx}
\makeindex

%--- термин: жирное начертание в тексте + метка + запись в указатель
\makeatletter
\newcommand{\term}[2]{\textbf{#1}\phantomsection\label{term:#2}\term@index}
\newcommand{\term@index}{\begingroup\@sanitize\term@index@}
\newcommand{\term@index@}[1]{\endgroup\index{#1}}
\makeatother

%--- сквозная нумерация сносок
\usepackage{chngcntr}
\counterwithout{footnote}{chapter}

%--- цвета
\definecolor{green_sea}{HTML}{2E8B57} % цвет рамки потока терминов
\definecolor{palette_green_9}{HTML}{E8F4E8} % «Зелёный 9» из палитры (COLORS.md)
\colorlet{table_header}{palette_green_9} % заливка шапки таблиц
\definecolor{methyl_orange_acid}{RGB}{235, 65, 30}
\definecolor{methyl_orange_base}{RGB}{255, 225, 0}
\definecolor{litmus_acid}{RGB}{205, 40, 65}
\definecolor{litmus_base}{RGB}{70, 90, 200}
\definecolor{phenolphthalein_base}{RGB}{255, 115, 175}
\definecolor{methyl_violet_I_acid}{RGB}{255, 235, 0}
\definecolor{methyl_violet_I_base}{RGB}{0, 150, 140}
\definecolor{methyl_violet_II_acid}{RGB}{0, 150, 140}
\definecolor{methyl_violet_II_base}{RGB}{20, 60, 180}
\definecolor{thymol_blue_I_acid}{RGB}{220, 30, 40}
\definecolor{thymol_blue_I_base}{RGB}{255, 225, 0}
\definecolor{methyl_violet_III_acid}{RGB}{20, 60, 180}
\definecolor{methyl_violet_III_base}{RGB}{120, 30, 150}
\definecolor{methyl_red_acid}{RGB}{215, 40, 40}
\definecolor{methyl_red_base}{RGB}{255, 220, 60}
\definecolor{bromthymol_blue_acid}{RGB}{220, 210, 40}
\definecolor{bromthymol_blue_base}{RGB}{30, 110, 180}
\definecolor{phenol_red_acid}{RGB}{240, 210, 60}
\definecolor{phenol_red_base}{RGB}{200, 30, 50}
\definecolor{thymol_blue_II_acid}{RGB}{255, 225, 0}
\definecolor{thymol_blue_II_base}{RGB}{30, 80, 170}
\definecolor{indigo_carmine_acid}{RGB}{30, 60, 140}
\definecolor{indigo_carmine_base}{RGB}{220, 210, 60}

%--- форматирование химических формул и уравнений
\usepackage[version=4]{mhchem}
\mhchemoptions{layout=stacked} % выравнивание заряда и количества атомов по вертикали

%--- создание окружений reaction и reaction*
\makeatletter
\newcounter{reaction}
\renewcommand\thereaction{C\,\thechapter.\arabic{reaction}}
\@addtoreset{reaction}{chapter}
\newcommand\reactiontag%
{\refstepcounter{reaction}\tag{\thereaction}}
\newcommand\reaction@[2][]%
{\begin{equation}\ce{#2}%
\ifx\@empty#1\@empty\else\label{#1}\fi%
\reactiontag\end{equation}}
\newcommand\reaction@nonumber[1]%
{\begin{equation*}\ce{#1}\end{equation*}}
\newcommand\reaction%
{\@ifstar{\reaction@nonumber}{\reaction@}}
\makeatother

%--- система форматирования чисел и единиц измерения
\usepackage{siunitx}
\usepackage{cancel}
\sisetup{
    output-decimal-marker = {,},
    group-minimum-digits = 5,
    group-digits = integer, % группировка только целой части числа
    exponent-product = \cdot,
    inter-unit-product = \cdot,
    per-mode = symbol, % группировка только целой части числа
    sticky-per
}

%--- масса
\DeclareSIUnit\ruKg{\text{кг}}
\DeclareSIUnit\ruT{\text{т}}
\DeclareSIUnit\ruG{\text{г}}
\DeclareSIUnit\ruDa{\text{Да}}

%--- длина
\DeclareSIUnit\ruKm{\text{км}}
\DeclareSIUnit\ruM{\text{м}}
\DeclareSIUnit\ruCm{\text{см}}
\DeclareSIUnit\ruMm{\text{мм}}
\DeclareSIUnit\ruUm{\text{мкм}}
\DeclareSIUnit\ruNm{\text{нм}}

%--- объём
\DeclareSIUnit\ruL{\text{л}}
\DeclareSIUnit\ruMl{\text{мл}}
\DeclareSIUnit\ruKl{\text{кл}}

%--- заряд
\DeclareSIUnit\ruC{\text{Кл}}
\DeclareSIUnit\ruEV{\text{эВ}}

%--- количество вещества
\DeclareSIUnit\ruMole{\text{моль}}
\DeclareSIUnit\ruKmole{\text{кмоль}}
\DeclareSIUnit\ruMmole{\text{ммоль}}

%--- давление
\DeclareSIUnit\ruPa{\text{Па}}
\DeclareSIUnit\ruKPa{\text{кПа}}
\DeclareSIUnit\ruMPa{\text{МПа}}
\DeclareSIUnit\ruAtm{\text{атм}}
\DeclareSIUnit\ruBar{\text{бар}}

%--- время
\DeclareSIUnit\ruS{\text{с}}
\DeclareSIUnit\ruMin{\text{мин}}
\DeclareSIUnit\ruCh{\text{ч}}

%--- энергия
\DeclareSIUnit\ruJ{\text{Дж}}
\DeclareSIUnit\ruKJ{\text{кДж}}
\DeclareSIUnit\ruMJ{\text{МДж}}

%--- мощность
\DeclareSIUnit\ruW{\text{Вт}}

%--- концентрация
\DeclareSIUnit\ruMolar{\text{М}}

%--- температура
\DeclareSIUnit\ruK{\text{К}}

%--- кубики для таблицы микросостояний
\newcommand{\diewidth}{6.1mm} % ширина одного кубика
\newcommand{\wdie}[1]{\includegraphics[width=\diewidth]{figs/entropy/die_w#1}}
\newcommand{\rdie}[1]{\includegraphics[width=\diewidth]{figs/entropy/die_r#1}}
\newcommand{\dpair}[2]{\mbox{\raisebox{-0.45\height}{\wdie{#1}\hspace{0.7mm}\rdie{#2}}}}

%--- hyperref
\usepackage{hyperref}
\hypersetup{
    pdftitle={Учебник по общей химии},
    pdfauthor={Михаил Коверда},
    pdfsubject={Учебник по общей химии для начинающих.
                  © 2026 Михаил Коверда. Лицензия CC BY-NC-SA 4.0
                  (кроме материалов в общественном достоянии):
                  https://creativecommons.org/licenses/by-nc-sa/4.0/},
    pdfkeywords={химия, общая химия, учебник, CC BY-NC-SA 4.0},
    pdflang={ru},
    pdfdisplaydoctitle=true,
    bookmarksnumbered=true,
}